\section{}\label{sec:two}

\subsection*{Стационарная теория возмущений}

\hspace{1em}В реальной квантовой механике часто приходится сталкиваться с системами, аналитическое решение которых не представляется возможным в связи с особенностями динамики этих систем. В таком случае обычно обращаются к аппроксимационным методам, используя которые можно найти решение для системы в необходимом приближении. Один из таких методов, который мы назвали квазиклассическим методом, был рассмотрен в последнем семинаре первого семестра. В нём нам важно было решить уравнение Шрёдингера для систем, который достаточно близки к классическим, хотя имеют чисто квантовые эффекты. Напомню, что в основу метода входит разложение функции, относительно которой решается дифференциальное уравнение, в ряд по малому параметру $\hbar$.

Попробуем развить новый аппроксимационный метод исходя из следующих соображений: пусть мы знаем, как описать систему в стационарном состоянии, то есть у нас есть решения уравнения
\[
\hat{H}^{(0)}\ket{\psi^{(0)}_n} = E^{(0)}_n\ket{\psi^{(0)}_n}
\]
Добавим к системе малые возмущения. Это может быть слабое электромагнитное поле, какой-то сдвиг системы или создания неровностей в потенциальной яме. В общем, любое достаточно малое возмущение. Тогда, обозначив это возмущение как $\lambda \hat{V}$, итоговый гамильтониан будет выглядеть следующим образом:
\[
\hat{H} = \hat{H}^{(0)} + \lambda\hat{V}
\]
Параметр $\lambda$ выберем маленьким, после чего сделаем разложение по нему в ряд для состояния и энергии:
\begin{gather*}
\ket{\psi} = \ket{\psi^{(0)}_n} + \ket{\psi^{(1)}_n} + \lambda^2\ket{\psi^{(2)}_n}+ \dots\\
E_n = E^{(0)}_n + \lambda E^{(1)}_n + \lambda^2 E^{(2)}_n + \dots
\end{gather*}

Метод, в котором считаются поправки для энергии и состояния возмущенного гамильтониана называется \textit{теория возмущений} (ТВ). В этом семинаре будет рассмотрена только стационарная ТВ, то есть с независящим от времени гамильтонианом. В первой части мы рассмотрим невырожденный случай (для каждого состояния только одно собственное значение), во второй части -- вырожденный.

\subsubsection*{Невырожденный случай}
Чтобы найти поправки разного порядка, подставим в исходное выражение полученные выше разложения:
\begin{multline*}
\left(H^{(0)} + \lambda V\right)\left(\ket{\psi^{(0)}_n} + \lambda\ket{\psi^{(1)}_n} +\lambda^2\ket{\psi^{(2)}_n} + \dots\right) =\\ = \left(E^{(0)}_n + \lambda E^{(1)}_n + \lambda^2E^{(2)}_n + \dots\right)\left(\ket{\psi^{(0)}_n} + \lambda\ket{\psi^{(1)}_n} +\lambda^2\ket{\psi^{(2)}_n} + \dots\right)
\end{multline*}
Соберём вместе члены с одинаковой степенью $\lambda$:
\begin{multline*}
    H^{(0)}\ket{\psi^{(0)}_n} + \lambda \left(H^{(0)}\ket{\psi^{(1)}_n} + V\ket{\psi^{(0)}_n}\right) + \lambda^2 \left(H^{(0)}\ket{\psi^{(2)}}_n + V\ket{\psi^{(1)}_n}\right) + \dots =\\= E^{(0)}_n\ket{\psi^{(0)}_n} + \lambda \left(E^{(0)}_n\ket{\psi^{(1)}_n} + E^{(1)}_n\ket{\psi^{(0)}_n}\right) + \lambda^2 \left(E^{(0)}_n\ket{\psi^{(2)}_n} + E^{(1)}_n\ket{\psi^{(1)}_n} + E^{(2)}_n\ket{\psi^{(0)}_n}\right) + \dots
\end{multline*}
Из этого уравнения лего заметить, что нулевая поправка (то есть члены с $\lambda^{0}$) есть в точности первоначальное уравнение с невозмущенным гамильтонианом.

Посчитаем первую поправку к энергии и к состоянию, то есть $E^{(1)}_n$ и $\ket{\psi^{(1)}_n}$. Для этого оставим члены с $\lambda$ в первой степени:
\[
H^{(0)}\ket{\psi^{(1)}_n} + V \ket{\psi^{(0)}_n} = E^{(0)}_n\ket{\psi^{(1)}_n} + E_n^{(1)}\ket{\psi^{(0)}_n}
\]
Чтобы получить первую поправку к энергии, домножим это выражение слева на $\bra{\psi^{(0)}_n}$:
\begin{gather*}
\bra{\psi^{(0)}_n}H^{(0)}\ket{\psi^{(1)}_n} + \bra{\psi^{(0)}_n}V\ket{\psi^{(0)}_n} = \bra{\psi^{(0)}_n}E^{(0)}_n\ket{\psi^{(1)}_n} + \bra{\psi^{(0)}_n}E^{(1)}_n\ket{\psi^{(0)}_n} \implies\\ 
\bra{\psi^{(0)}_n}E^{(0)}_n\ket{\psi^{(1)}_n} + \bra{\psi^{(0)}_n}V\ket{\psi^{(0)}_n} = \bra{\psi^{(0)}_n}E^{(0)}_n\ket{\psi^{(1)}_n} + E^{(1)}_n \implies\\
E^{(1)}_n = \bra{\psi^{(0)}_n}V\ket{\psi^{(0)}_n}
\end{gather*}
Получается, что первая поправка по энергии это среднее значение возмущения в невозмущенном состоянии. Этот результат является очень важным и имеет большое количество практических приложений, на которые мы обязательно посмотрим в этом семестре.

Далее, найдём первую поправку для состояния. Разделим эту процедуру на две части: сначала найдём поправку для индексов $m \neq n$, затем обсудим ситуацию с $m = n$. Выразим первую поправку через нулевую как 
\[
\ket{\psi^{(1)}_n} = \sum\limits_{m\neq n}c_m^{(1)}\ket{\psi^{(0)}_m}
\]
Это можно сделать, так как невозмущенное состояние образует полный базис в пространстве, значит через него может быть выражено любое другое состояние.

Чтобы получить поправки к энергии, мы домножили уравнение на состояние с тем же индексом $\bra{\psi^{(0)}_n}$, получив таким образом диагональные члены гамильтониана. Сейчас мы поступим иначе, домножив уравнение на $\bra{\psi^{(0)}_m}$, где $m\neq n$:
\begin{gather*}
    \bra{\psi^{(0)}_m}\sum\limits_{m\neq n}c^{(1)}_mH^{(0)}\ket{\psi^{(0)}_m} + \bra{\psi^{(0)}_m}V\ket{\psi^{(0)}_n} = \bra{\psi^{(0)}_m}\sum\limits_{m\neq n}c^{(1)}_mE^{(0)}_n\ket{\psi^{(0)}_m} + \bra{\psi^{(0)}_m}E^{(1)}_n\ket{\psi^{(0)}_n}\implies\\
    c^{(1)}_m E^{(0)}_m + \bra{\psi^{(0)}_m}V\ket{\psi^{(0)}_n} = c^{(1)}_m E^{(0)}_n \implies\\
    c^{(1)}_m  =  \frac{\bra{\psi^{(0)}_m}V\ket{\psi^{(0)}_n}}{\left(E^{(0)}_n - E^{(0)}_m\right)}
\end{gather*}
Итак, коэффициенты для первой поправки состояния есть недиагональные матричные элементы с поправкой на расстояние между энергетическими уровнями системы. Теперь рассмотри ситуацию, когда $m=n$. Напрямую достать значение коэффициента $c^{(1)}_n$ из уравнения выше нельзя, поэтому мы воспользуемся условием нормировки. Пусть $\braket{\psi}{\psi} = 1$ с точностью до порядка поправок, то есть $\sum_{l=0}^{p}\braket{\psi^{(l)}_n}{\psi^{(p-l)}_n} = 0$. Тогда для первой поправки $\braket{\psi^{(1)}_n}{\psi^{(0)}_n} = \braket{\psi^{(0)}_n}{\psi^{(1)}_n}$. Это условие выполняется при $c^{(1)}_n = 0$.

Имея эти выражения для первых поправок, мы уже можем решать достаточно широкий класс задач, в которых встречается малое возмущение. Удовлетворение условию малости доказывается апостериорно, то есть сначала применяется ТВ, а затем рассматриваются условия $E^{(p)}_n \ll E^{(p-1)}_n$ и $\braket{\psi^{(p)}_n}{\psi^{(p)}_n} \ll \braket{\psi^{(p-1)}_n}{\psi^{(p-1)}_n}$.

Однако, нередко бывают случаи, когда первый порядок зануляется. В таком случае необходимо посчитать поправку второго порядка. Давайте сделаем это, начав с поправки к энергии. Для этого проделаем те же самые шаги, то есть подставим нужное количество элементов и домножим слева на $\bra{\psi^{(0)}_n}$:
\begin{gather*}
    \bra{\psi^{(0)}_n}H^{(0)}\ket{\psi^{(2)}_n} + \bra{\psi^{(0)}_n}V\ket{\psi^{(1)}_n} = \bra{\psi^{(0)}_n}E^{(0)}_n\ket{\psi^{(2)}_n} + \bra{\psi^{(0)}_n}E^{(1)}_n\ket{\psi^{(1)}_n} + E^{(2)}_n \implies\\ 
    E^{(2)}_n = \bra{\psi^{(0)}_n}V\ket{\psi^{(1)}_n} -E^{(1)}_n\bra{\psi^{(0)}_n}\ket{\psi^{(1)}_n} 
\end{gather*}
Из предыдущих уравнений мы знаем, что $\braket{\psi^{(0)}_n}{\psi^{(1)}_n}=\sum_{m\neq n} c^{(1)}_{m} \braket{\psi^{(0)}_n}{\psi^{(0)}_m}$. Для второго члена справа это выражение даст ноль. Если же мы разложим первый член таким образом, то получим связь второй поправки для энергии с коэффициентами для первой поправки:
\[
 E^{(2)}_n = \sum\limits_{m\neq n}c^{(1)}_{m}\bra{\psi^{(0)}_n}V\ket{\psi^{(0)}_m} = \sum\limits_{m\neq n}\frac{|\bra{\psi^{(0)}_n}V\ket{\psi^{(0)}_m}|^2}{(E^{(0)}_n - E^{(0)}_m)}
\]

Для вычисления второй поправки проделаем тот же самый алгоритм -- разложим вторую поправку к состоянию по нулевой:
\[
\ket{\psi^{(2)}_n} = \sum\limits_{m\neq n}c^{(2)}_m\ket{\psi^{(0)}_m}
\]
Теперь домножим уравнение слева на $\bra{\psi^{(0)}_m}$:
\begin{gather*}
     \bra{\psi^{(0)}_m}\sum\limits_{m\neq n}c^{(2)}_mH^{(0)}\ket{\psi^{(0)}_m} + \bra{\psi^{(0)}_m}V\ket{\psi^{(1)}_n} = \bra{\psi^{(0)}_m}\sum\limits_{m\neq n}c^{(2)}_mE^{(0)}_n\ket{\psi^{(0)}_m} + \\
     + \bra{\psi^{(0)}_m}E^{(1)}_n\ket{\psi^{(1)}_n} + \bra{\psi^{(0)}_m}E^{(2)}_n\ket{\psi^{(0)}_n} \implies \\
     c^{(2)}_mE^{(0)}_m + \bra{\psi^{(0)}_m}\sum\limits_{k\neq n}c^{(1)}_k V\ket{\psi^{(0)}_k} = c^{(2)}_mE^{(0)}_n\ + \bra{\psi^{(0)}_m}\sum\limits_{k\neq n}c^{(1)}_kE^{(1)}_n\ket{\psi^{(0)}_k}\implies \\
     c^{(2)}_{m} = \sum\limits_{k\neq n} \frac{\bra{\psi^{(0)}_m}V\ket{\psi^{(0)}_k} \bra{\psi^{(0)}_k}V\ket{\psi^{(0)}_n}} {(E^{(0)}_n - E^{(0)}_k)(E^{(0)}_k - E^{(0)}_m)}-\frac{\bra{\psi^{(0)}_m}V\ket{\psi^{(0)}_n} \bra{\psi^{(0)}_n}V\ket{\psi^{(0)}_n}}{(E^{(0)}_n - E^{(0)}_m)^2}
\end{gather*}

Рассмотрим ситуацию, когда $m=n$. Тогда, исходя из нормировки, указанной выше, получим
\begin{gather*}
\braket{\psi^{(2)}_n}{\psi^{(0)}_n} + \braket{\psi^{(1)}_n}{\psi^{(1)}_n} + \braket{\psi^{(0)}_n}{\psi^{(2)}_n} = 0 \implies\\
c^{(2)}_n = -\frac{1}{2}\sum\limits_{m\neq n}\frac{|\bra{\psi^{(0)}_m}V\ket{\psi^{(0)}_n}|^2}{(E^{(0)}_n - E^{(0)}_m)^2}
\end{gather*}

Любую поправку можно посчитать подобным образом. Но, конечно, так как алгоритм идентичный, существует итеративные формулы для получения любой поправки:
\begin{gather*}
    E^{(p)}_n = \sum_{k}\bra{\psi^{(0)}_n}V\ket{\psi^{(0)}_k}c^{(p-1)} - \sum_{l=1}^{p-1}E^{(p-l)}_nc^{(l)}\\
    c^{(p)}_{m\neq n} = \frac{1}{E^{(0)}_{n} - E^{(0)}_{m}}\left(\sum_{k}\bra{\psi^{(0)}_m}V\ket{\psi^{(0)}_k}c^{(p-1)} - \sum_{l=0}^{p-1}E^{(p-l)}_m c^{(l)}\right) \\
    c^{(p)}_{n} = -\frac{1}{2}\sum\limits_{l=1}^{p-1}\sum_{k}(c^{(l)}_k)^*c^{(p-l)_k}
\end{gather*}

Итак, у нас есть все необходимые формулы, чтобы решить первую задачу на стационарную невырожденную задачу.
\excersize
\begin{center}
\textit{Найти с помощью теории возмущений поправки к уровням энергии линейного гармонического осциллятора, возмущенного полем вида:}
\[
a)\,\hat{V} = \alpha \hat{x},\quad b)\,\hat{V} = A\hat{x}^3 + B\hat{x}^4
\]
\end{center}
Ещё раз выпишем формулы, которыми мы будем пользоваться:
\[
E_n^{(1)} = \bra{\psi_n^{(0)}}V\ket{\psi_n^{(0)}}, \quad E_n^{(2)} = \sum\limits_{m\neq n} \frac{\left| \bra{\psi_n^{(0)}} V \ket{\psi_m^{(0)}} \right|^2}{E_n^{(0)}-E_m^{(0)}}.
\]
Для начала нам нужно получить решение невозмущённой задачи. В 6 семинаре прошлого семестра мы подробно обсуждали гармонический осциллятор, поэтому сейчас можем просто выписать решения:
\[
E_n^{(0)} = \hbar\omega\left(n+\frac{1}{2}\right), \quad \ket{n} = \frac{(\hat{a}^{\dagger})^n}{\sqrt{n!}}\ket{0}
\]
Так как мы пользуемся записью состояний через оператор рождения, запишем в таком же представлении и оператор координаты: $\hat{x} = x_0(\hat{a}+\hat{a}^{\dagger})/{\sqrt{2}}$, $x_0 = \sqrt{\hbar/(m\omega)}$.

Тогда, можем сразу посчитать первую поправку по энергии в случае $\hat{V}=\alpha \hat{x}$:
\[
E_n^{(1)} = \alpha\frac{x_0}{\sqrt{2}}\left(\bra{n}\hat{a}\ket{n}+\bra{n}\hat{a}^{\dagger}\ket{n}\right) = \alpha\frac{x_0}{\sqrt{2}}\left(\sqrt{n}\braket{n}{n-1}+\sqrt{n+1}\braket{n}{n+1}\right) = 0
\]

Теперь можем приступить к вычислению второй поправки. Для удобства, сначала отдельно посчитаем недиагональные матричные элементы оператора возмущения:
\[
    \bra{\psi_n^{(0)}}\hat{V}\ket{\psi_m^{(0)}} = \frac{\alpha x_0}{\sqrt{2}}\left(\bra{n}\hat{a}\ket{m}+\bra{n}\hat{a}^{\dagger}\ket{m}\right) = \frac{\alpha x_0}{\sqrt{2}}\left(\sqrt{m}\delta_{n(m-1)}+\sqrt{m+1}\delta_{n(m+1)}\right)
\]
Подставим их в выражение для вторых поправок:
\begin{multline*}
    E_n^{(2)} = \frac{\alpha^2 x_0^2}{2} \sum\limits_{m\neq n} \frac{\left|\sqrt{m} \delta_{n(m-1)} +\sqrt{m+1}\delta_{n(m+1)}\right|^2}{\hbar\omega(n-m)} =\\= \frac{\alpha^2 x_0^2}{2} \sum\limits_{m\neq n} \left(\frac{m\delta_{n(m-1)}}{\hbar\omega(n-m)}+\frac{(m+1)\delta_{n(m+1)}}{\hbar\omega(n-m)}\right) = \frac{\alpha^2 x_0^2}{2}\left(-\frac{n+1}{\hbar\omega}+\frac{n}{\hbar\omega}\right) = -\frac{\alpha^2}{2m\omega^2}
\end{multline*}

Посмотрим, что происходит в случае $\hat{V} = A\hat{x}^3+B\hat{x}^4$. В 20 упражнении 6 семинара мы уже считали средние значения операторов $\hat{x}^3$ и $\hat{x}^4$, поэтому сейчас снова можем воспользоваться готовыми результатами и выписать первую поправку:
\[
E_n^{(1)} = \bra{n}A\hat{x}^3\ket{n} + \bra{n}B\hat{x}^4\ket{n} = 0 + B\frac{x_0^4}{4}(6n^2+6n+3) = \frac{3}{2}Bx_0^4\left(n^2+n+\frac{1}{2}\right)
\]

Со второй поправкой дела обстоят несколько сложнее. Для её вычисления, нам нужно честно посчитать следующие матричные элементы: $\bra{n}\hat{x}^4\ket{m}$ и $\bra{n}\hat{x}^3\ket{m}$. Идейно это сделать несложно, можно просто  ``в лоб'' действовать этими операторами на состояние $\ket{m}$, а затем посчитать скалярное произведение (не забывая, что мы рассматриваем только случаи $m\neq n$). Здесь же приведём сразу посчитанные результаты:
\begin{multline*}
    \bra{m}\frac{x_0^3}{2\sqrt{2}}(\hat{a}+\hat{a}^{\dagger})\ket{n} = \frac{x_0^3}{2\sqrt{2}}\left(\sqrt{m(m-1)(m-2)}\delta_{n(m-3)}+3m^{3/2}\delta_{n(m-1)}+\right. \\ \left.+ 3(m+1)^{3/2}\delta_{n(m+1)}+\sqrt{(m+1)(m+2)(m+3)}\delta_{n(m+3)}\right)
\end{multline*}
\begin{multline*}
    \bra{m}\frac{x_0^4}{4}(\hat{a}+\hat{a}^{\dagger})\ket{n} = \frac{x_0^4}{4}\left(\sqrt{m(m-1)(m-2)(m-3)}\delta_{n(m-4)}+(4m-2)\sqrt{m(m-1)}\delta_{n(m-2)}+ \right.\\ \left.+(4m+6)\sqrt{(m+1)(m+2)}\delta_{n(m+2)}+\sqrt{(m+1)(m+2)(m+3)(m+4)}\delta_{n(m+4)}\right)
\end{multline*}
Просуммируем полученные элементы по всем состояниям:
\begin{align*}
    \sum\limits_{m\neq n}\frac{\left|\bra{n}\hat{x}^3\ket{m}\right|^2}{\hbar\omega(n-m)} &= \frac{x_0^6}{8\hbar\omega}\left(-\frac{1}{3}n(n-1)(n-2)-9n^3+9(n+1)^3+\frac{1}{3}(n+1)(n+2)(n+3)\right)=\\
    &=\frac{30n^2+30n+11}{8}\frac{x_0^6}{\hbar\omega} \\
    \sum\limits_{m\neq n}\frac{\left|\bra{n}\hat{x}^4\ket{m}\right|^2}{\hbar\omega(n-m)}&=\frac{x_0^8}{16\hbar\omega}\left(-\frac{1}{4}n(n-1)(n-2)(n-3)-2(2n-1)^2n(n-1)\right.+ \\
    &\left.+2(2n+3)^2(n+1)(n+2)+\frac{1}{4}(n+1)(n+2)(n+3)(n+4)\right) = \\
    &= \frac{34n^3+51n^2+59n+21}{8}\frac{x_0^8}{\hbar\omega}
\end{align*}

Тогда, наша поправка выглядит так:
\[
E_n^{(2)} = A^2\frac{34n^3+51n^2+59n+21}{8}\frac{x_0^8}{\hbar\omega}+B^2\frac{30n^2+30n+11}{8}\frac{x_0^6}{\hbar\omega}
\]
\csquare{darklavender}

\subsubsection*{Вырожденный случай}
В реальной жизни часто приходится сталкиваться с вырожденными системами (например, энергетические уровни в атомах). Напомню, что система называется вырожденной, если одной и той же энергии соответствуют несколько состояний. Теория, описанная выше, не подходит для описания вырожденных систем, так как она предполагает, что при уменьшении возмущения ($\lambda \rightarrow 0$) конкретное состояние приходит к невозмущенной энергии $E^{(0)}_n$. В вырожденном случае любая линейная комбинация из вырожденных состояний имеет одинаковую невозмущенную энергию, поэтому не понятно, к какой конкретной линейной комбинации придёт возмущённые состояния при уменьшении возмущения. 

Чтобы решить эту неоднозначность, необходимо проделать алгоритм, аналогичный невозмущенной теории, но сразу подставлять линейную комбинацию для невозмущенного состояния $\ket{\psi^{(0)}_n} = \sum_{\alpha}c^{(0)}_{n, \alpha}\ket{\psi^{(0)}_{n,\alpha}}$. Далее греческими индексами будет выражаться вырождение, а латинскими, как и раньше, номера уровней.

Посчитаем таким образом первый порядок возмущения для вырожденной системы, домножив уравнение для первого порядка по $\lambda$ на $\bra{\psi^{(0)}_{n, \beta}}$:
\begin{gather*}
    \hat{H}_0\ket{\psi^{(1)}_{n}} + \hat{V}\ket{\psi^{(0)}_{n}} = E^{(0)}_n\ket{\psi^{(1)}_{n}} + E^{(1)}_n\ket{\psi^{(0)}_{n}} \bigg{|} \bra{\psi^{(0)}_{n, \beta}}\; \cdot \implies \\
    E^{(0)}_n\braket{\psi^{(0)}_{n, \beta}}{\psi^{(1)}_{n}} + \sum\limits_{\alpha}c^{(0)}_{n, \alpha} \bra{\psi^{(0)}_{n, \beta}} \hat{V} \ket{\psi^{(0)}_{n, \alpha}} = E^{(0)}_n\braket{\psi^{(0)}_{n, \beta}}{\psi^{(1)}_{n}} + E^{(1)}_n\sum\limits_{\alpha}c^{(0)}_{n, \alpha}\braket{\psi^{(0)}_{n, \beta}}{\psi^{(0)}_{n, \alpha}} \implies \\
    \sum\limits_{\alpha}c^{(0)}_{n, \alpha} \bra{\psi^{(0)}_{n, \beta}} \hat{V} \ket{\psi^{(0)}_{n, \alpha}} = E^{(1)}_nc^{(0)}_{n, \beta} \implies \\
    \sum\limits_{\alpha}\left(\bra{\psi^{(0)}_{n, \beta}} \hat{V} \ket{\psi^{(0)}_{n, \alpha}} - E^{(1)}_n \delta_{\beta\alpha}\right)c^{(0)}_{n, \alpha} = 0
\end{gather*}

Это уравнение -- в точности стационарное уравнение Шрёдингера, с тем отличием, что вместо гамильтониана слева стоит матрица $N \times N$, где $N$ -- порядок вырождения. Решить эту задачу сильно проще, чем решать задачу для целого гамильтониана. Хотя решается она также: для нахождения первой поправки к энергии нам необходимо решить \textit{секулярное уравнение}
\[
\det(\hat{V} - E^{(1)}_n) = 0.
\]
Полученные таким образом собственные значения и будут первыми поправками к энергии. Если эти собственные значения будут различны, то говорят о \textit{снятии вырождения}, то есть состояние перешло из вырожденного (с одинаковыми энергиями) к невырожденному (с разными энергиями). Собственные состояния, соответствующие найденным собственным значениям, называются \textit{правильными}.

Теперь получим первую поправку к вырожденному состоянию. Для этого, как и до этого, домножим уравнение первого порядка на $\bra{\psi^{(0)}_{m}}, \; m\neq n$, перед этим представив первую поправку как:
\[
\ket{\psi^{(1)}_{m}} =\sum\limits_{k\neq n}c^{(1)}_{k}\ket{\psi^{(0)}_{k}}
\]
Здесь нет никакого вырождения, так как мы предположили, что вырожден только уровень $n$. Тогда, проецируя уравнение на $\bra{\psi^{(0)}_{m}}$, получим
\begin{gather*}
    \hat{H}_0\ket{\psi^{(1)}_{n}} + \hat{V}\ket{\psi^{(0)}_{n}} = E^{(0)}_n\ket{\psi^{(1)}_{n}} + E^{(1)}_n\ket{\psi^{(0)}_{n}} \bigg{|} \bra{\psi^{(0)}_{m}}\; \cdot \implies\\
    \bra{\psi^{(0)}_{m}}\sum\limits_{k}c^{(1)}_kE^{(0)}_k\ket{\psi^{(0)}_{k}} +  \bra{\psi^{(0)}_{m}}\sum\limits_{\alpha}c^{(0)}_{n, \alpha}\hat{V} \ket{\psi^{(0)}_{n, \alpha}} = E^{(0)}_n \bra{\psi^{(0)}_{m}}\sum\limits_{k}c^{(1)}_k\ket{\psi^{(0)}_{k}} \implies\\
    c^{(1)}_m = \sum\limits_{\alpha}c^{(0)}_{n, \alpha}\frac{\bra{\psi^{(0)}_m} \hat{V} \ket{\psi^{(0)}_{n, \alpha}}}{E^{(0)}_{n} - E^{(0)}_m} = \frac{\bra{\psi^{(0)}_m} \hat{V} \ket{\psi^{(0)}_{n}}}{E^{(0)}_{n} - E^{(0)}_m}
\end{gather*}
Итак, для первой поправки состояний разницы между вырожденным и невырожденным случаем нет. Учитывая нормировку, можно установить, что $c^{(1)}_{n, \alpha} = 0$. Для подсчёта вторых поправок семинар и так уже затянулся, поэтому эта задача предлагается для рассмотрения читателям самостоятельно, так как подход не изменится. Здесь же рассмотрим задачу, в которой явно видно практическую пользу теории возмущения.
\excersize
\begin{center}
\textit{Найти смещение уровня энергии основного состояния атома водорода под влиянием конечных размеров ядра. Ядро считать равномерно заряженным по объему шаром радиуса $r_0$}
\end{center}

Для начала разберёмся, как повлияет конечный размер ядра на гамильтониан. Обычно мы считаем, что ядро является точечным зарядом, то есть $U(r) = -e^2/r$ на всем пространстве. В случае же равномерно заряженного по объему шара известного радиуса, потенциал будет следующим:
\[
U(r)=\begin{cases}
    -\frac{e^2}{r}, &r\in [r_0, +\infty) \\
    \frac{e^2}{2r_0}\left(\frac{r^2}{r_0^2}-3\right), &r\in(0, r_0)
\end{cases}
\]
Кратко обсудим то, откуда взялся этот потенциал. В общем случае потенциал в точке $\vec{r}$ с известным распределением заряда задается формулой $\varphi(\vec{r}) = \int_{\mathbb{R}^3}\rho(\vec{r})/|\vec{r}-\vec{r}'|d^3r$, где $\rho(\vec{r}')$ -- плотность заряда, распределенного внутри заряженного шара, $\vec{r}$ -- радиус-вектор точки наблюдения, $\vec{r}'$ -- радиус-вектор точки внутри заряженного шара. Тогда, считая заряд распределенным равномерно по всему объему ядра, мы можем найти плотность: $\rho(r)=3Q/4\pi r_0^3$. Далее, подставляя это значение в интеграл и переходя к сферическим координатам (так как в рамках данной задачи явно прослеживается сферическая симметрия), мы и получаем нашу итоговую формулу. Тогда, посчитаем наше возмущение:
\[
V(r) = \begin{cases}
    0, & r\in[r_0, +\infty) \\
    \frac{e^2}{2r_0}\left(\frac{r^2}{r_0^2}-3+\frac{2r_0}{r}\right), & r\in(0, r_0)
\end{cases}
\]
Видим, что оно отличается от нуля только внутри самого ядра. При этом, мы знаем характерный размер ядра водорода $r_0\sim 10^{-15}$м, а средний радиус орбиты электрона (боровский радиус) -- $a\sim 10^{-11}$м. То есть, возмущение, вызванное конечными размерами ядра водорода, действует на очень далеком расстоянии от электрона. Поэтому мы можем считать это возмущение малым и применить к нему теорию возмущений.

Для того, чтобы посчитать смещение уровня энергии основного состояния атома водорода, нам нужно выписать волновую функцию этого состояния (мы подробно обсуждали ее вывод в $9$ семинаре прошлого семестра):
\[
\psi_{100}(r, \theta, \varphi) = \sqrt{\frac{1}{\pi a^3}}e^{-\frac{r}{a}},
\]
где $a$ -- радиус Бора.

Наконец, можем вычислить поправку по энергии, считая $r_0 \ll a$, то есть $e^{-\frac{r}{a}}\approx 1$, где $r\in(0, r_0)$:
\begin{align*}
    E_{100}^{(1)} &= \bra{\psi_{100}^{(0)}}\hat{V}\ket{\psi_{100}^{(0)}} = \frac{1}{\pi a^3} \frac{e^2}{2r_0} \int\limits_{0}^{2\pi}d\varphi\int\limits_{0}^{\pi}\sin\theta d\theta\int\limits_{0}^{r_0}r^2e^{-\frac{2r}{a}}\left(\frac{r^2}{r_0}-3+\frac{2r_0}{r}\right)dr = \\
    &= \frac{2e^2}{a^3r_0}\int\limits_0^{r_0}\left(\frac{r^4}{r_0}-3r^2+2rr_0\right)dr = \frac{2e^2}{a^3}\frac{r_0^2}{5} = -\frac{4}{5}\left(\frac{r_0}{a}\right)^2E_0^{(0)}
\end{align*}
Получается, поправка составляет примерно $10^{-8}$ часть от исходной энергии, что еще раз подтверждает, что теория возмущений в данной ситуации применима
\csquare{darklavender}

\excersize
\begin{center}
\textit{Найти расщепление уровня энергии атома водорода с $n = 2$ в однородном электрическом поле (эффект Штарка). Найти правильные волновые функции нулевого приближения}
\end{center}
В невозмущенном случае при $n=2$ кратность вырождения равна $g=n^2=4$, то есть у нас имеется $4$ волновые функции, соответствующие одному значению энергии $E_2^{(0)} = -\frac{1}{2n^2}\frac{me^4}{\hbar^2}$. Выпишем эти функции:
\begin{align*}
    \psi_{200}(r,\theta,\varphi)&=\sqrt{\frac{1}{8\pi a^3}}\left(1-\frac{r}{2a}\right)e^{-\frac{r}{2a}}\\
    \psi_{210}(r,\theta,\varphi)&=\sqrt{\frac{1}{32\pi a^5}}re^{-\frac{r}{2a}}\cos\theta \\
    \psi_{21\pm1}(r,\theta,\varphi)&=\sqrt{\frac{1}{64\pi a^5}}re^{-\frac{r}{2a}}\sin\theta e^{\pm i\varphi},
\end{align*}
где $a$ -- Боровский радиус.

Гамильтониан невозмущенной задачи нам давно известен: $\hat{H}_0 = \frac{\hat{p}^2}{2m}-\frac{e^2}{r}$. Добавляя однородное электрическое поле, мы получаем гамильтониан возмущенной задачи: $\hat{H} = \hat{H}_0 - e(\vec{r}, \vec{\mathcal{E}})$, где $\vec{\mathcal{E}}$ -- напряженность нашего внешнего поля. Для удобства будем считать, что $\vec{\mathcal{E}}=(0, 0, \mathcal{E})$, то есть $\hat{V} = -e\mathcal{E}\hat{z}$.

Для нахождения правильных волновых функций нам сначала надо найти матрицу $\bra{\psi^{(0)}_{n, \beta}} \hat{V} \ket{\psi^{(0)}_{n, \alpha}}$. В $9$ семинаре прошлого семестра мы решали очень похожую задачу (упражнение $29$). Подробный подсчет можно посмотреть там, а мы же просто воспользуемся результатами: обнулятся все матричные элементы, кроме $\bra{\psi^{(0)}_{200}} \hat{V} \ket{\psi^{(0)}_{210}}$ и симметричного ему $\bra{\psi^{(0)}_{210}} \hat{V} \ket{\psi^{(0)}_{200}}$:
\[
\bra{\psi^{(0)}_{200}} \hat{V} \ket{\psi^{(0)}_{210}}=\bra{\psi^{(0)}_{210}} \hat{V} \ket{\psi^{(0)}_{200}}=3e\mathcal{E}a
\]
Следовательно, матричное представление оператора возмущения (для удобства будем считать, что $\psi_1=\psi_{211}, \, \psi_2=\psi_{210}, \, \psi_3=\psi_{21-1}, \, \psi_4=\psi_{200}$:
\[V_{\alpha\beta}=
\begin{pmatrix}
    0 & 0 & 0 & 0 \\ 0 & 0 & 0 & 3e\mathcal{E}a \\  0 & 0 & 0 & 0 \\ 0 & 3e\mathcal{E}a & 0 & 0
\end{pmatrix}
\]
Наконец, можем решить секулярное уравнение $\det(V_{\alpha\beta}-E_n^{(1)}\delta_{\alpha\beta}) = 0$:
\[
    \begin{vmatrix}
        -E_n^{(1)} & 0 & 0 & 0 \\ 0 & -E_n^{(1)} & 0 & 3e\mathcal{E}a \\  0 & 0 & -E_n^{(1)} & 0 \\ 0 & 3e\mathcal{E}a & 0 & -E_n^{(1)}
    \end{vmatrix}=E_n^{(1)4}-9e^2\mathcal{E}^2a^2E_n^{(1)2} =0 \implies \left[\begin{matrix}\
        E_n^{(1)2} = 0 \\ E_n^{(1)} = \pm3e\mathcal{E}a 
    \end{matrix}\right.
\]
Итого, поправки первого порядка к уровню энергии $E_2^{(0)}$:
\[
E_{1,4}^{(1)}=\mp3e\mathcal{E}a, \quad E_{2,3}^{(1)}=0
\]
Найдем правильные волновые функции. Они представляются в виде линейной комбинации волновых функций невозмущенной задачи:
\[
\Tilde{\psi} = c_1\psi_1+c_2\psi_2+c_3\psi_3+c_4\psi_4
\]
с неизвестными пока коэффициентами $c_i$. Для их нахождения решим следующую систему уравнений:
\[
\begin{pmatrix}
     -E_n^{(1)} & 0 & 0 & 0 \\ 0 & -E_n^{(1)} & 0 & 3e\mathcal{E}a \\  0 & 0 & -E_n^{(1)} & 0 \\ 0 & 3e\mathcal{E}a & 0 & -E_n^{(1)}
\end{pmatrix}
\begin{pmatrix}
    c_1 \\ c_2 \\ c_3 \\ c_4
\end{pmatrix}=0 \implies\begin{cases}
    E_n^{(1)}c_1 = E_n^{(1)}c_3 = 0 \\
    E_n^{(1)} с_2 + 3e\mathcal{E} a c_4=0 \\
    3e\mathcal{E}a c_2 - E_n^{(1)} c_4 =0 \\
    |c_1|^2+|c_2|^2+|c_3|^2+|c_4|^2=1
\end{cases}
\]
Последнее уравнение является условием нормировки. Подставим в данную систему все поправки, которые мы посчитали:
\begin{enumerate}
    \item $E^{(1)}_{2,3}=0$. В таком случае $c_2=c_4=0$, а $|c_1|^2+|c_3|^2=1$. Для $E^{(1)}_2$ можем взять $c_1=1, \, c_3 = 0, \, \Tilde{\psi}_1 = \psi_{211}$. Аналогично, для $E^{(1)}_3$ подходит $c_3=1, \, c_1=0, \, \Tilde{\psi}_3=\psi_{21-1}$.
    \item $E_1=-3e\mathcal{E}a$. Отсюда следует, что $c_1=c_3=0$, $c_2=-1/\sqrt{2},\, c_4 = 1/\sqrt{2}$. Правильная волновая функция выглядит следующим образом: $\Tilde{\psi}_1 = 1/\sqrt{2}(\psi_{200}-\psi_{210})$.
    \item $E_4=3e\mathcal{E}a$. Аналогично предыдущему случаю, $c_1=c_3=0$, $c_2=c_4=1/\sqrt{2}$, $\Tilde{\psi}_4 = 1/\sqrt{2}(\psi_{200}+\psi_{210})$.
\end{enumerate}
Видно, что постоянное электрическое поле сняло вырождение, но не полностью, так как уровень $E_{2,3}^{(1)}$ двукратно вырожден.
\csquare{darklavender}

В следующем семинаре поговорим о нестационарной теории возмущения.